\section{Introduction}
\label{sec:intro}

\subsection*{The Saint--Venant inequality and its quantitative form}

Let $\Omega\subset\R^2$ be a bounded open set of finite perimeter with
$|\Omega|=\pi$.  The \emph{torsion function} of $\Omega$ is the unique
solution $u\in H^1_0(\Omega)$ of $-\Delta u=1$ in $\Omega$, $u=0$ on
$\partial\Omega$; its integral $T(\Omega):=\int_\Omega u$ is the
\emph{torsional rigidity}.  The classical Saint--Venant conjecture,
proved by P\'olya and Szeg\H{o} \cite{PolyaSzego1951} via the
Schwarz symmetrization of the torsion potential, asserts that the disc
maximises torsional rigidity at fixed area:
\[
  T(\Omega)\ \le\ T(B_1)\ =\ \frac{\pi}{8},
\]
with equality if and only if $\Omega$ is a disc.  A more precise, and
substantially more difficult, question is to quantify stability: how
large must $\Omega$ deviate from a disc when the deficit
$\dT(\Omega):=\tfrac\pi8-T(\Omega)$ is small?  The natural measure of
deviation is the \emph{Fraenkel asymmetry}
\[
  \AF(\Omega):=\min_{x_0\in\R^2}
    \frac{|\Omega\triangle B(x_0,r_\Omega)|}{|\Omega|},\qquad
  r_\Omega:=\sqrt{|\Omega|/\pi},
\]
which is zero if and only if $\Omega$ is a disc.  The sharp quantitative
Saint--Venant inequality is the statement
\begin{equation}\label{eq:sharp-main}
  \AF(\Omega)^2\ \le\ C\,\dT(\Omega)
\end{equation}
for some \emph{absolute} constant $C$ (independent of $\Omega$, of $|\Omega|$,
and of any regularity of $\partial\Omega$).  The exponent $2$ on the
left and the exponent $1$ on the right are both optimal: the
ratio $\AF(\Omega)^2/\dT(\Omega)$ is bounded and bounded away from
zero for a one-parameter family of ellipses.

Talenti \cite{Talenti1976} proved a pointwise comparison between the
decreasing rearrangement of $u_\Omega$ and the torsion function of the
equimeasure disc, giving $T(\Omega)\le T(B_1)$ and the quantitative
comparison $\mu(v)\le\mu_B(v)$ between the distribution functions.
This Talenti comparison, which we state precisely as \cref{thm:talenti},
is a key input in the present work.  Talenti's result also implies the
profile identity $\int_0^{1/4}(\mu_B-\mu)\d v=\dT$ (see \cref{lem:profile}),
which localizes the deficit to a single integral of a non-negative
integrand.

\subsection*{Prior work on quantitative stability}

The Faber--Krahn problem for the first Dirichlet eigenvalue $\lambda_1$
is the structural analogue of the Saint--Venant problem, and the two
share much of their quantitative theory.  For the quantitative
isoperimetric inequality the sharp form
$\AF(E)\le C_{\mathrm{F}}\,\diso(E)^{1/2}$ (where $\diso(E)$ denotes the
isoperimetric deficit and $C_{\mathrm F}$ is an absolute constant; see
\cref{thm:fmp}) was established by Fusco, Maggi, and Pratelli
\cite{FMP2008}; this result, which we use as a black box, is the
\emph{only} external quantitative inequality the present paper imports.

For Saint--Venant/Faber--Krahn stability, Brasco and Pratelli
\cite{BrascoPratelli2012} proved a quantitative stability result for
both the torsional rigidity and the first eigenvalue, establishing
\eqref{eq:sharp-main} with a \emph{non-sharp} exponent.  The sharp
inequality \eqref{eq:sharp-main} --- with exponent $2$ and an absolute
constant, valid with no regularity on $\partial\Omega$ --- was
established by Brasco, De Philippis, and Velichkov \cite{BDV2015}.
Their argument proceeds by a \emph{selection principle}: from a
near-extremal sequence one extracts a subsequence that, after
appropriate normalisation, converges in a strong topology to a
$C^{1,\gamma}$ domain; the stability of the latter is then proved by a
second-variation analysis.  The selection principle is a powerful and
elegant tool, but it is non-constructive, non-quantitative in the
convergence, and it relies on the regularity theory for minimizers of
shape functionals in a fundamental way.  A survey of quantitative
spectral inequalities, including the state of the art up to~2017, is
given by Brasco and De Philippis \cite{BrascoDePhilippis2017}.

\subsection*{The main theorem}

The main result of this paper is the following, proved in \cref{sec:main}
as \cref{thm:main}.

\begin{theorem}[Main theorem; see \cref{thm:main}]
  There exists an absolute constant $C$ such that every bounded open set
  $\Omega\subset\R^2$ with $|\Omega|=\pi$ satisfies
  \[
    \AF(\Omega)^2\ \le\ C\,\dT(\Omega).
  \]
\end{theorem}

This is the same conclusion as \cite{BDV2015}, but obtained by an
entirely different route.  The proof requires no selection principle,
no mass transport, and no symmetry-type argument (Serrin overde\-termined
problem / Magnanini--Poggesi estimates).  The only external inputs are:
classical function-space and geometric-measure-theory tools
(the coarea formula, the planar isoperimetric inequality, the Talenti
comparison, and the sharp quantitative isoperimetric inequality of
\cite{FMP2008}), all quoted with precise references in
\cref{sec:prelim}.

\subsection*{Strategy and what is new}

The argument replaces the selection--regularity--second-variation
architecture of \cite{BDV2015} by a purely level-set construction
driven by the two \emph{level deficits}:
\begin{align*}
  D_H(v) &:= \mu(v)\bigl(-\mu'(v)\bigr) - P(\{u>v\})^2\ \ge\ 0,\\
  D_I(v) &:= P(\{u>v\})^2 - 4\pi\mu(v)\ \ge\ 0,
\end{align*}
where $\mu(v)=|\{u>v\}|$ is the distribution function and
$P(\{u>v\})$ is the perimeter of the superlevel set.  The central
identity (\cref{thm:level-identity}, proved in \cref{sec:identity})
\[
  \int_0^m\bigl(D_H(v)+D_I(v)\bigr)\d v\ =\ 4\pi\,\dT(\Omega)
\]
encodes the entire deficit in a non-negative integral over the level
parameter.  All the estimates in the paper are consequences of this
identity and the classical Talenti comparison.

\smallskip
\noindent\textbf{The two sharp ingredients.}
Two novel observations make the reduction exact:

\begin{enumerate}
\item \emph{Fold-content bound with no square-root loss.}
  At the optimal inset level $\hat v\sim\sqrt{\dT}$, the connected
  component $\Etil$ of $\{u>\hat v\}$ of largest measure is trapped in a
  thin annulus $B_{\rho_i}(z)\subset\Etil\subset B_{\rho_e}(z)$ of
  width $\eta=\rho_e-\rho_i\lesssim\dT^{1/8}$ (established in
  \cref{sec:annulus}).  The total angular measure of rays along which
  $\Etil$ is \emph{not} a radial graph is bounded, via the
  point-source flux identity and the cancellation
  $|I_{\Etil}-2\pi\hat r|\le a\eta/(\rho_i\hat r)$
  (where $a=|\Etil\triangle B_{\hat r}(z)|$), by a quantity of order
  $\Exc+a\eta$ — without the factor of $\sqrt{\diso}$ that the
  naive application of the quantitative isoperimetric inequality
  \cite{FMP2008} would introduce.  This is the content of
  \cref{prop:fold,cor:gd} in \cref{sec:fold}.

\item \emph{Cutting caps decreases the deficit.}
  To manufacture a genuine radial graph $G$ from $\Etil$ one must excise
  the ``caps'' (the portions of $\Etil$ that cause each ray to intersect
  $\partial\Etil$ more than once).  A potential concern is that this
  excision increases the scale-invariant deficit
  $\dstar(G):=(T(B_G)-T(G))/T(B_G)$.  In fact it does not: the
  equal-area comparison ball shrinks faster than torsion is lost, so
  $\dstar$ \emph{decreases}.  This is \cref{lem:cutdeficit} in
  \cref{sec:cut}, proved via the Green-function representation of
  torsional rigidity and a barrier estimate in the collar.
\end{enumerate}

With $\tau:=\sqrt{\dT}$, the two observations combine to produce a
continuous radial graph $G$ with amplitude $\lesssim\dT^{1/8}$,
scale-invariant deficit $\dstar(G)\le C\dstar(\Omega)$, and symmetric
difference $|\Omega\triangle G|\le C\sqrt{\dT}$ (the precise reduction
is \cref{thm:reduction} in \cref{sec:main}).  The
perturbative stability of nearly spherical sets
(\cref{thm:graphG} in \cref{sec:nearly}) then gives $\AF(G)^2\le C_G\dstar(G)$,
and the triangle inequality for the symmetric difference transfers this
to $\AF(\Omega)^2\le C\dT(\Omega)$.

\subsection*{Organization of the paper}

\Cref{sec:prelim} collects the classical prerequisites: the theory of
sets of finite perimeter (De Giorgi--Federer structure theorem, coarea
formula, differentiation of the distribution function), the planar
isoperimetric inequality and the sharp quantitative version
\cite{FMP2008}, the interior regularity and analyticity of the torsion
function \cite{GilbargTrudinger2001}, and Sard's theorem.

\Cref{sec:identity} introduces the two level deficits $D_H$ and $D_I$,
proves their non-negativity, establishes the central level-set deficit
identity $\int_0^m(D_H+D_I)\d v=4\pi\dT$ (\cref{thm:level-identity}),
and derives the Talenti comparison and the profile identity.

\Cref{sec:inset} carries out the level-set construction: it proves
deficit transfer to any superlevel set, selects the good level $\hat v$
by an averaging argument, establishes the closeness of $\{u>\hat v\}$
to $\Omega$, and isolates the inset $\Etil$ as the dominant connected
component.

\Cref{sec:annulus} proves that a connected hole-free finite-perimeter
set with controlled isoperimetric deficit is trapped in a thin
concentric annulus, using the angular-perimeter coarea identity and a
crossing-number argument.

\Cref{sec:fold} bounds the fold content of $\Etil$ --- the obstruction
to being a radial graph --- via the point-source flux identities and the
thin-annulus cancellation, achieving the sharp estimate without the
quantitative-isoperimetric square-root loss.

\Cref{sec:cut} proves that excising the caps to manufacture a radial
graph decreases the scale-invariant torsion deficit, via a
Green-function comparison and a barrier estimate.

\Cref{sec:nearly} establishes the perturbative stability of nearly
spherical sets for the torsion functional, proving that any radial graph
of small amplitude satisfies $\AF^2\le C_G\dstar$, with explicit
absolute constants, by a second-order expansion of the torsion deficit
in Fourier modes.

\Cref{sec:main} assembles all pieces, optimises the level parameter at
$\tau=\sqrt{\dT}$, and derives the main theorem in full.

\subsection*{Notation}

Throughout the paper, $\Omega\subset\R^2$ is a bounded open set with
$|\Omega|=\pi$.  The unit disc is $B_1$; $B(x_0,r)$ denotes the open
disc of radius $r$ centred at $x_0$.  The symbol $C$ denotes a generic
positive absolute constant whose value may change from line to line;
named constants such as $C_G$, $\eta_G$, $A_0$, $\delta_*$ are fixed
at their first occurrence.  All \emph{absolute} constants are
independent of $\Omega$ and of any geometric or regularity data of
$\partial\Omega$.  The notation $f\lesssim g$ means $f\le Cg$ for an
absolute $C$.
